\documentclass[aspectratio=169]{beamer}

% Theme and color scheme
\usetheme{default}
\usecolortheme{default}

% Remove navigation symbols
\setbeamertemplate{navigation symbols}{}

% Simple color scheme for printing
\definecolor{darkblue}{RGB}{25,25,112}
\setbeamercolor{title}{fg=darkblue}
\setbeamercolor{frametitle}{fg=darkblue}
\setbeamercolor{structure}{fg=darkblue}

% Simple footline
\setbeamertemplate{footline}[frame number]

% Packages
\usepackage{amsmath}
\usepackage{amsfonts}
\usepackage{amssymb}
\usepackage{graphicx}
\usepackage{tikz}
\usepackage{booktabs}
\usepackage{array}
\usepackage{listings}
\usepackage{xcolor}

% Set graphics path for images (relative to slides directory)
\graphicspath{{../Images/}}

% Title information
\title[Chapter 3.13: Harvard Architecture]{Structured Computer Architecture\\Chapter 3.13: Harvard Architecture\\Fifth-Order Systems (5-OS)}
\author{Based on Structured Computer Architecture: An Order-Based Approach}
\date{}

% Custom command for highlighting concepts
\newcommand{\concept}[1]{\textbf{\color{darkblue}#1}}

\begin{document}

% Title slide
\begin{frame}
\titlepage
\end{frame}

% Reset section counter for this chapter
\setcounter{section}{0}

% Define sections for this chapter
\section{Historical Context and Motivation}
\section{Interpreting vs. Executing}
\section{Harvard Computer Architecture}
\section{RISC Approach}
\section{Memory Hierarchy and Performance}
\section{Implementation Considerations}

% Table of contents
\begin{frame}{Outline}
\tableofcontents
\end{frame}

% Learning objectives
\begin{frame}{Learning Objectives}
By the end of this chapter, you will understand:
\begin{itemize}
\item The historical development from Turing machines to Harvard architecture
\item The distinction between interpreting and executing instructions
\item Harvard vs. Von Neumann architectural models
\item RISC (Reduced Instruction Set Computer) principles
\item Memory hierarchy implications in Harvard architecture
\item Performance benefits of separate instruction and data memories
\item Fifth-order system characteristics and loop structures
\end{itemize}
\end{frame}

% Section 1: Historical Context and Motivation
\section{Historical Context and Motivation}

\begin{frame}{From Turing Machine to Harvard Architecture}
\begin{columns}
\begin{column}{0.6\textwidth}
\concept{Universal Turing Machine (UTM) Insight}:
\begin{itemize}
\item Turing divided tape into two parts:
\begin{itemize}
\item Machine description $M(T)$ (program)
\item Associated tape $T$ (data)
\end{itemize}
\item Two-memory concept: programs and data
\item Multi-tape TM: multiple memories, multiple heads
\item Same-cycle access to different memories
\end{itemize}

\vspace{0.3cm}
\concept{Mathematical Foundation}:
\begin{itemize}
\item UTM operates with dual memory structure
\item Multiple read heads enable parallel access
\item Theoretical basis for Harvard model
\end{itemize}
\end{column}
\begin{column}{0.4\textwidth}
\begin{center}
\textbf{UTM Memory Structure}
\vspace{0.3cm}

\tikz{
% Program memory
\draw[thick] (0,2.5) rectangle (2,3);
\node at (1,2.75) {Program $M(T)$};

% Data memory  
\draw[thick] (0,1.5) rectangle (2,2);
\node at (1,1.75) {Data $T$};

% Processing unit
\draw[thick] (0.5,0.5) rectangle (1.5,1);
\node at (1,0.75) {UTM};

% Connections
\draw[<->] (1,1) -- (1,1.5);
\draw[<->] (1,2) -- (1,2.5);
}
\end{center}
\end{column}
\end{columns}
\end{frame}

\begin{frame}{Historical Milestone: Harvard Mark I (1944)}
\begin{columns}
\begin{column}{0.6\textwidth}
\concept{Key Historical Event}:
\begin{itemize}
\item \textbf{March 29, 1944}: John von Neumann runs program on Harvard Mark I
\item Working on Manhattan Project
\item First practical implementation of dual-memory concept
\end{itemize}

\vspace{0.3cm}
\concept{Howard Aiken's Design Decision}:
\begin{itemize}
\item Architect chose two-memory organization
\item One memory for programs
\item Separate memory for data
\item Inspired by Turing Machine theory
\end{itemize}

\vspace{0.3cm}
\concept{Birth of Harvard Abstract Model}:
\begin{itemize}
\item Parallel development to Von Neumann model
\item Both inspired by Turing Machine concepts
\item Different approaches to memory organization
\end{itemize}
\end{column}
\begin{column}{0.4\textwidth}
\begin{center}
\textbf{Harvard Mark I Concept}
\vspace{0.3cm}

\tikz{
% Timeline
\draw[thick] (0,0) -- (0,3);
\node at (-0.5,3) {1944};
\node at (-0.5,2) {Mark I};
\node at (-0.5,1) {Aiken};
\node at (-0.5,0) {UTM};

% Architecture
\draw[thick] (1,2.5) rectangle (2.5,3);
\node at (1.75,2.75) {Program};
\draw[thick] (1,1.5) rectangle (2.5,2);
\node at (1.75,1.75) {Data};
\draw[thick] (1.25,0.5) rectangle (2.25,1);
\node at (1.75,0.75) {CPU};

\draw[->] (1.75,1) -- (1.75,1.5);
\draw[->] (1.75,2) -- (1.75,2.5);
}
\end{center}
\end{column}
\end{columns}
\end{frame}

% Section 2: Interpreting vs. Executing
\section{Interpreting vs. Executing}

\begin{frame}{UTM Processing Collaboration}
\begin{columns}
\begin{column}{0.6\textwidth}
\concept{UTM Processing Model}:
\begin{itemize}
\item Collaboration between program and finite automaton
\item Each data change requires multiple steps:
\begin{enumerate}
\item Reading from program part of tape
\item State transitions of UTM's finite automaton
\end{enumerate}
\end{itemize}

\vspace{0.3cm}
\concept{Computational Complexity Factors}:
\begin{enumerate}
\item \textbf{Program size}: Flexible, redesignable structure
\item \textbf{Automaton size}: Fixed, physically implemented structure
\end{enumerate}

\vspace{0.3cm}
\concept{Design Considerations}:
\begin{itemize}
\item Hardware complexity reduction is crucial
\item Specification, design, verification, validation, testing effort
\item Minimize control automaton complexity
\end{itemize}
\end{column}
\begin{column}{0.4\textwidth}
\begin{center}
\textbf{UTM Processing Flow}
\vspace{0.3cm}

\tikz{
% Program tape
\draw[thick] (0,3) rectangle (2,3.5);
\node at (1,3.25) {Program};

% Data tape
\draw[thick] (0,2) rectangle (2,2.5);
\node at (1,2.25) {Data};

% Finite automaton
\draw[thick] (0.5,0.5) rectangle (1.5,1.5);
\node at (1,1) {FA};

% Process flow
\draw[->] (1,2.5) -- (1,1.5);
\draw[->] (1,3) -- (1,1.5);
\node at (2.2,1.75) {Read};
\node at (2.2,1.25) {Transition};
}
\end{center}
\end{column}
\end{columns}
\end{frame}

\begin{frame}{Shannon's Contribution and Zero-State Solution}
\begin{columns}
\begin{column}{0.6\textwidth}
\concept{Claude Shannon's Approach (1956)}:
\begin{itemize}
\item Reduced finite automaton to only 2 states
\item Consequence: Increased binary circuit complexity
\item State reduction ≠ complexity reduction
\item Larger next-state calculation circuits
\end{itemize}

\vspace{0.3cm}
\concept{Zero-State Finite Automaton}:
\begin{itemize}
\item Ultimate solution: eliminate finite automaton entirely
\item Zero internal states = zero control complexity
\item Justifies complete elimination of control complexity
\item Direct instruction execution without interpretation
\end{itemize}

\vspace{0.3cm}
\concept{Design Philosophy}:
\begin{itemize}
\item Minimize hardware control complexity
\item Prefer execution over interpretation
\item UTM with zero states as ideal model
\end{itemize}
\end{column}
\begin{column}{0.4\textwidth}
\begin{center}
\textbf{Control Complexity Evolution}
\vspace{0.3cm}

\tikz{
% Traditional FA
\draw[thick] (0,3) circle (0.3);
\draw[thick] (0.7,3) circle (0.3);
\draw[thick] (1.4,3) circle (0.3);
\node at (0.7,2.5) {Many States};

% Shannon's 2-state
\draw[thick] (0,2) circle (0.3);
\draw[thick] (0.7,2) circle (0.3);
\node at (0.35,1.5) {2 States};
\node at (0.35,1.2) {Complex Logic};

% Zero-state
\draw[thick] (0.35,0.5) rectangle (0.7,0.8);
\node at (0.525,0.65) {Direct};
\node at (0.525,0.2) {Zero States};
}
\end{center}
\end{column}
\end{columns}
\end{frame}

\begin{frame}{Interpreting vs. Executing Instructions}
\begin{columns}
\begin{column}{0.6\textwidth}
\concept{Two Approaches to Instruction Processing}:

\vspace{0.3cm}
\textbf{1. Interpretation}:
\begin{itemize}
\item Instruction is \textbf{interpreted}
\item Executes sequence of micro-instructions
\item Generated by control automaton
\item Multiple clock cycles per instruction
\item Complex control logic required
\end{itemize}

\vspace{0.3cm}
\textbf{2. Execution}:
\begin{itemize}
\item Instruction is \textbf{executed} directly
\item Single clock cycle execution
\item Instruction fields applied directly to components
\item Minimal control complexity
\item UTM with zero internal states model
\end{itemize}

\vspace{0.3cm}
\concept{Executive Structure Preference}:
\begin{itemize}
\item UTM with two tapes (Harvard model)
\item Zero internal states (direct execution)
\end{itemize}
\end{column}
\begin{column}{0.4\textwidth}
\begin{center}
\textbf{Instruction Processing}
\vspace{0.3cm}

\tikz{
% Interpretation path
\node at (1,3.5) {\textbf{Interpretation}};
\draw[thick] (0,3) rectangle (2,3.3);
\node at (1,3.15) {Instruction};
\draw[->] (1,3) -- (1,2.7);
\draw[thick] (0,2.4) rectangle (2,2.7);
\node at (1,2.55) {Control Automaton};
\draw[->] (1,2.4) -- (1,2.1);
\draw[thick] (0,1.8) rectangle (2,2.1);
\node at (1,1.95) {Micro-instructions};

% Execution path
\node at (1,1.3) {\textbf{Execution}};
\draw[thick] (0,0.8) rectangle (2,1.1);
\node at (1,0.95) {Instruction};
\draw[->] (1,0.8) -- (1,0.5);
\draw[thick] (0,0.2) rectangle (2,0.5);
\node at (1,0.35) {Direct to Hardware};
}
\end{center>
\end{column>
\end{columns>
\end{frame>

% Section 3: Harvard Computer Architecture
\section{Harvard Computer Architecture}

\begin{frame}{Harvard Abstract Model: Fifth-Order System}
\begin{columns}
\begin{column}{0.6\textwidth}
\concept{Fifth Loop Structure}:
\begin{itemize}
\item Processor + Program Memory = 4-OS (Fourth-Order System)
\item Adding Data Memory loop = 5-OS (Fifth-Order System)
\item Dual memory connections enable parallel access
\end{itemize}

\vspace{0.3cm}
\concept{Key Architectural Features}:
\begin{itemize}
\item Separate program and data memories
\item Independent memory buses
\item Simultaneous instruction and data access
\item Enhanced memory bandwidth
\item Reduced memory access conflicts
\end{itemize}

\vspace{0.3cm}
\concept{RISC Implementation Focus}:
\begin{itemize}
\item Harvard model naturally suits RISC processors
\item Simple instruction formats
\item Direct execution capability
\item Minimal control complexity
\end{itemize}
\end{column}
\begin{column}{0.4\textwidth}
\begin{center}
\textbf{Harvard 5-OS Model}
\vspace{0.3cm}

\tikz{
% Program Memory
\draw[thick] (0,3) rectangle (1.5,3.5);
\node at (0.75,3.25) {Program};
\node at (0.75,3.05) {Memory};

% Data Memory
\draw[thick] (0,1) rectangle (1.5,1.5);
\node at (0.75,1.25) {Data};
\node at (0.75,1.05) {Memory};

% Processor
\draw[thick] (2,2) rectangle (3,2.5);
\node at (2.5,2.25) {CPU};

% Connections (dual bus)
\draw[<->] (1.5,3.25) -- (2,2.4);
\draw[<->] (1.5,1.25) -- (2,2.1);

% Loop indicators
\node at (1.75,2.8) {4-OS};
\node at (1.75,1.7) {5-OS};
}
\end{center}
\end{column}
\end{columns}
\end{frame}

\begin{frame}{Harvard vs. Von Neumann: Key Differences}
\begin{columns}
\begin{column}{0.6\textwidth}
\concept{Harvard Architecture Advantages}:
\begin{itemize}
\item \textbf{Dual connections}: Independent program and data paths
\item \textbf{Parallel access}: Instruction fetch + data access in same cycle
\item \textbf{No memory conflicts}: Separate buses eliminate contention
\item \textbf{Higher bandwidth}: Two memory accesses per cycle possible
\end{itemize>

\vspace{0.3cm}
\concept{Memory Hierarchy Impact}:
\begin{itemize}
\item Deeper memory hierarchy possible
\item L1 cache miss resolution much faster
\item L2, L3 cache levels more effective
\item Severe consequences for memory access speed optimization
\end{itemize}

\vspace{0.3cm}
\concept{RALU Subsystem Benefits}:
\begin{itemize}
\item Simultaneous instruction and data processing
\item Reduced pipeline stalls
\item Enhanced throughput
\end{itemize}
\end{column}
\begin{column}{0.4\textwidth}
\begin{center}
\textbf{Architecture Comparison}
\vspace{0.3cm}

\tikz{
% Von Neumann
\node at (1,3.5) {\textbf{Von Neumann}};
\draw[thick] (0,3) rectangle (2,3.3);
\node at (1,3.15) {Unified Memory};
\draw[thick] (0.5,2.5) rectangle (1.5,2.8);
\node at (1,2.65) {CPU};
\draw[<->] (1,2.8) -- (1,3);

% Harvard
\node at (1,2) {\textbf{Harvard}};
\draw[thick] (0,1.5) rectangle (1,1.7);
\node at (0.5,1.6) {Prog};
\draw[thick] (1.2,1.5) rectangle (2.2,1.7);
\node at (1.7,1.6) {Data};
\draw[thick] (0.6,1) rectangle (1.6,1.3);
\node at (1.1,1.15) {CPU};
\draw[<->] (0.8,1.3) -- (0.5,1.5);
\draw[<->] (1.4,1.3) -- (1.7,1.5);
}
\end{center}
\end{column}
\end{columns}
\end{frame}

% Section 4: RISC Approach
\section{RISC Approach}

\begin{frame}{Reduced Instruction Set Computer (RISC) Principles}
\begin{columns}
\begin{column}{0.6\textwidth}
\concept{RISC Philosophy}:
\begin{itemize}
\item Maintain only simplest, most frequently used instructions
\item Large instruction set → reduced instruction set
\item Rejected instructions implemented as sequences
\item All complex operations decomposed into simple ones
\end{itemize}

\vspace{0.3cm}
\concept{Performance Advantages}:
\begin{itemize}
\item Simple design outperforms complex CPU designs
\item Huge amounts of performance data support RISC
\item Machine code only marginally larger than CISC
\item Better performance despite larger code size
\end{itemize>

\vspace{0.3cm}
\concept{Design Benefits}:
\begin{itemize}
\item Simplified control logic
\item Faster instruction execution
\item Easier pipeline implementation
\item Reduced hardware complexity
\item Better compiler optimization opportunities
\end{itemize}
\end{column}
\begin{column}{0.4\textwidth}
\begin{center}
\textbf{RISC vs CISC}
\vspace{0.3cm}

\tikz{
% CISC
\node at (1,3.5) {\textbf{CISC}};
\draw[thick] (0,3) rectangle (2,3.3);
\node at (1,3.15) {Complex Instructions};
\draw[thick] (0,2.7) rectangle (2,3);
\node at (1,2.85) {Complex Control};
\node at (1,2.5) {Slower, Smaller Code};

% RISC
\node at (1,2) {\textbf{RISC}};
\draw[thick] (0,1.5) rectangle (2,1.7);
\node at (1,1.6) {Simple Instructions};
\draw[thick] (0,1.3) rectangle (2,1.5);
\node at (1,1.4) {Simple Control};
\node at (1,1) {Faster, Larger Code};

% Performance arrow
\draw[->, thick, red] (2.2,2.75) -- (2.2,1.55);
\node[red] at (2.7,2.15) {Better};
\node[red] at (2.7,1.95) {Performance};
}
\end{center}
\end{column}
\end{columns}
\end{frame}

\begin{frame}{RISC Instruction Characteristics}
\begin{columns}
\begin{column}{0.6\textwidth}
\concept{Instruction Set Properties}:
\begin{itemize}
\item \textbf{Fixed instruction length}: Simplifies decode logic
\item \textbf{Simple addressing modes}: Reduces complexity
\item \textbf{Load/store architecture}: Memory access only via specific instructions
\item \textbf{Large register file}: Reduces memory accesses
\end{itemize}

\vspace{0.3cm}
\concept{Execution Model}:
\begin{itemize}
\item One instruction per clock cycle (ideally)
\item Direct execution without micro-code
\item Minimal control unit complexity
\item Pipeline-friendly design
\end{itemize}

\vspace{0.3cm}
\concept{Harvard Architecture Synergy}:
\begin{itemize}
\item RISC naturally fits Harvard model
\item Separate instruction and data paths
\item Simultaneous fetch and execute
\item Optimal for pipeline implementation
\end{itemize}
\end{column}
\begin{column}{0.4\textwidth}
\begin{center}
\textbf{RISC Instruction Format}
\vspace{0.3cm}

\tikz{
% Fixed format instruction
\draw[thick] (0,2.5) rectangle (3,3);
\draw[thick] (0.6,2.5) -- (0.6,3);
\draw[thick] (1.2,2.5) -- (1.2,3);
\draw[thick] (1.8,2.5) -- (1.8,3);
\draw[thick] (2.4,2.5) -- (2.4,3);

\node at (0.3,2.75) {Op};
\node at (0.9,2.75) {Rs};
\node at (1.5,2.75) {Rt};
\node at (2.1,2.75) {Rd};
\node at (2.7,2.75) {Func};

\node at (1.5,2.2) {32-bit Fixed Format};

% Simple operations
\node at (1.5,1.7) {\textbf{Simple Operations}};
\node at (1.5,1.4) {ADD, SUB, AND, OR};
\node at (1.5,1.1) {LOAD, STORE};
\node at (1.5,0.8) {BRANCH, JUMP};
}
\end{center>
\end{column>
\end{columns>
\end{frame>

% Section 5: Memory Hierarchy and Performance
\section{Memory Hierarchy and Performance}

\begin{frame}{Memory Hierarchy in Harvard Architecture}
\begin{columns}
\begin{column}{0.6\textwidth}
\concept{Hierarchical Memory Structure}:
\begin{itemize}
\item \textbf{L1 Cache}: Separate instruction and data caches
\item \textbf{L2 Cache}: Unified or separate second-level cache
\item \textbf{L3 Cache}: Third-level cache for larger systems
\item \textbf{Main Memory}: Separate or unified main memory
\end{itemize}

\vspace{0.3cm}
\concept{Performance Benefits}:
\begin{itemize}
\item L1 cache miss resolved much faster
\item Independent cache hierarchies reduce conflicts
\item Parallel access to instruction and data caches
\item Better cache utilization patterns
\end{itemize}

\vspace{0.3cm}
\concept{Access Patterns}:
\begin{itemize}
\item Instruction access: Sequential, predictable
\item Data access: Random, less predictable
\item Separate optimization strategies possible
\item Different cache policies for each type
\end{itemize>
\end{column>
\begin{column}{0.4\textwidth>
\begin{center>
\textbf{Memory Hierarchy}
\vspace{0.3cm}

\tikz{
% L1 Caches
\draw[thick] (0,3.5) rectangle (1,4);
\node at (0.5,3.75) {L1-I};
\draw[thick] (1.5,3.5) rectangle (2.5,4);
\node at (2,3.75) {L1-D};

% L2 Cache
\draw[thick] (0.5,2.8) rectangle (2,3.2);
\node at (1.25,3) {L2 Cache};

% L3 Cache
\draw[thick] (0.25,2.1) rectangle (2.25,2.5);
\node at (1.25,2.3) {L3 Cache};

% Main Memory
\draw[thick] (0,1.4) rectangle (2.5,1.8);
\node at (1.25,1.6) {Main Memory};

% CPU
\draw[thick] (0.75,4.2) rectangle (1.75,4.6);
\node at (1.25,4.4) {CPU};

% Connections
\draw[<->] (1,4.2) -- (0.5,4);
\draw[<->] (1.5,4.2) -- (2,4);
\draw[<->] (1.25,3.5) -- (1.25,3.2);
\draw[<->] (1.25,2.8) -- (1.25,2.5);
\draw[<->] (1.25,2.1) -- (1.25,1.8);
}
\end{center>
\end{column>
\end{columns>
\end{frame>

\begin{frame}{Performance Analysis and Optimization}
\begin{columns}
\begin{column}{0.6\textwidth}
\concept{Performance Metrics}:
\begin{itemize}
\item \textbf{Memory bandwidth}: Dual access capability
\item \textbf{Cache hit rates}: Separate optimization for I-cache and D-cache
\item \textbf{Pipeline efficiency}: Reduced stalls from memory conflicts
\item \textbf{Throughput}: Instructions per cycle improvement
\end{itemize}

\vspace{0.3cm}
\concept{Optimization Strategies}:
\begin{itemize}
\item Instruction cache: Optimize for sequential access
\item Data cache: Optimize for spatial/temporal locality
\item Prefetching: Different strategies for each cache
\item Cache sizing: Independent optimization
\end{itemize}

\vspace{0.3cm}
\concept{System-Level Benefits}:
\begin{itemize}
\item Reduced memory bottlenecks
\item Better resource utilization
\item Improved scalability
\item Enhanced parallel processing capability
\end{itemize>
\end{column>
\begin{column}{0.4\textwidth>
\begin{center>
\textbf{Performance Comparison}
\vspace{0.3cm}

\tikz{
% Performance bars
\draw[thick, fill=blue!30] (0,3) rectangle (1,3.5);
\node at (0.5,2.7) {Von Neumann};
\draw[thick, fill=green!30] (0,2) rectangle (1.8,2.5);
\node at (0.9,1.7) {Harvard};

\node at (1,3.8) {\textbf{Memory Bandwidth}};

% Cache performance
\draw[thick, fill=blue!30] (0,1) rectangle (1.2,1.3);
\node at (0.6,0.8) {Unified Cache};
\draw[thick, fill=green!30] (0,0.3) rectangle (1.6,0.6);
\node at (0.8,0.1) {Separate Caches};

\node at (1,1.5) {\textbf{Cache Efficiency}};
}
\end{center>
\end{column>
\end{columns>
\end{frame>

% Section 6: Implementation Considerations
\section{Implementation Considerations}

\begin{frame}{Hardware Implementation Challenges}
\begin{columns}
\begin{column}{0.6\textwidth}
\concept{Design Complexity}:
\begin{itemize}
\item \textbf{Dual memory interfaces}: Separate controllers needed
\item \textbf{Bus arbitration}: Managing multiple memory accesses
\item \textbf{Cache coherency}: Maintaining consistency between caches
\item \textbf{Address space management}: Separate or unified addressing
\end{itemize}

\vspace{0.3cm}
\concept{Cost Considerations}:
\begin{itemize}
\item Additional memory controllers
\item Separate cache structures
\item More complex interconnect
\item Higher pin count requirements
\end{itemize}

\vspace{0.3cm}
\concept{Design Trade-offs}:
\begin{itemize}
\item Performance vs. complexity
\item Cost vs. benefit analysis
\item Power consumption implications
\item Silicon area requirements
\end{itemize>
\end{column>
\begin{column}{0.4\textwidth>
\begin{center>
\textbf{Implementation Complexity}
\vspace{0.3cm}

\tikz{
% Complexity factors
\node at (1,3.5) {\textbf{Design Factors}};

\draw[thick] (0,3) rectangle (2,3.2);
\node at (1,3.1) {Memory Controllers};

\draw[thick] (0,2.6) rectangle (2,2.8);
\node at (1,2.7) {Cache Management};

\draw[thick] (0,2.2) rectangle (2,2.4);
\node at (1,2.3) {Bus Arbitration};

\draw[thick] (0,1.8) rectangle (2,2);
\node at (1,1.9) {Address Mapping};

% Trade-off arrow
\draw[->, thick] (0.5,1.5) -- (1.5,1.5);
\node at (0.3,1.3) {Cost};
\node at (1.7,1.3) {Performance};
}
\end{center>
\end{column>
\end{columns>
\end{frame>

\begin{frame}{Modern Applications and Variations}
\begin{columns}
\begin{column}{0.6\textwidth}
\concept{Contemporary Implementations}:
\begin{itemize}
\item \textbf{Digital Signal Processors (DSPs)}: Pure Harvard architecture
\item \textbf{Microcontrollers}: Modified Harvard with unified addressing
\item \textbf{ARM Processors}: Harvard pipeline with Von Neumann interface
\item \textbf{RISC-V}: Configurable memory architecture options
\end{itemize}

\vspace{0.3cm}
\concept{Hybrid Approaches}:
\begin{itemize}
\item Harvard pipeline with Von Neumann memory
\item Separate L1 caches, unified L2 cache
\item Dynamic switching between modes
\item Software-configurable memory organization
\end{itemize}

\vspace{0.3cm}
\concept{Future Directions}:
\begin{itemize}
\item Near-memory computing
\item Processing-in-memory architectures
\item Heterogeneous memory systems
\item AI-optimized memory hierarchies
\end{itemize>
\end{column>
\begin{column}{0.4\textwidth>
\begin{center>
\textbf{Modern Implementations}
\vspace{0.3cm}

\tikz{
% DSP
\draw[thick] (0,3.5) rectangle (2,3.8);
\node at (1,3.65) {DSP - Pure Harvard};

% MCU
\draw[thick] (0,3) rectangle (2,3.3);
\node at (1,3.15) {MCU - Modified};

% ARM
\draw[thick] (0,2.5) rectangle (2,2.8);
\node at (1,2.65) {ARM - Hybrid};

% RISC-V
\draw[thick] (0,2) rectangle (2,2.3);
\node at (1,2.15) {RISC-V - Configurable};

% Future
\draw[thick, dashed] (0,1.5) rectangle (2,1.8);
\node at (1,1.65) {Future - PIM};
}
\end{center>
\end{column>
\end{columns>
\end{frame>

% Summary
\section{Summary}

\begin{frame}{Chapter Summary}
\begin{columns}
\begin{column}{0.6\textwidth}
\concept{Key Concepts Covered}:
\begin{itemize}
\item Harvard architecture emerged from UTM dual-tape concept
\item Historical significance of Harvard Mark I (1944)
\item Distinction between instruction interpretation and execution
\item RISC principles naturally align with Harvard architecture
\item Fifth-order systems provide enhanced memory bandwidth
\item Performance benefits through separate instruction and data paths
\item Modern implementations use hybrid approaches
\end{itemize>

\vspace{0.3cm}
\concept{Architectural Advantages}:
\begin{itemize}
\item Parallel instruction and data access
\item Reduced memory conflicts
\item Enhanced pipeline efficiency
\item Better cache utilization
\item Improved overall system performance
\end{itemize>
\end{column>
\begin{column}{0.4\textwidth>
\begin{center>
\textbf{Harvard Architecture Legacy}
\vspace{0.3cm}

\tikz{
\node[draw, circle, minimum width=2cm] (center) at (0,0) {Harvard};
\node[draw, rectangle] (risc) at (0,1.5) {RISC};
\node[draw, rectangle] (perf) at (-1.5,0.5) {Performance};
\node[draw, rectangle] (mem) at (1.5,0.5) {Memory};
\node[draw, rectangle] (modern) at (0,-1.5) {Modern CPUs};

\draw[->] (center) -- (risc);
\draw[->] (center) -- (perf);
\draw[->] (center) -- (mem);
\draw[->] (center) -- (modern);
}
\end{center>
\end{column>
\end{columns>
\end{frame>

\begin{frame}{Discussion Questions}
\begin{enumerate}
\item How does the Harvard architecture address the Von Neumann bottleneck?
\item What are the trade-offs between Harvard and Von Neumann architectures?
\item Why do RISC processors naturally fit the Harvard model?
\item How does separate instruction and data memory affect cache design?
\item What are the implications of zero-state finite automaton concept?
\item How do modern processors implement hybrid Harvard-Von Neumann approaches?
\item What role does the memory hierarchy play in Harvard architecture performance?
\item How might future computing paradigms evolve the Harvard concept?
\end{enumerate}
\end{frame}

\end{document}

